\subsection{Koeficientët e reflektimit dhe transmetimit}
Kur një valë harmonike bie normalisht mbi një ndërfaqe, fusha në anën e parë shkruhet si shumë e valës rënëse dhe të reflektuar,
\begin{equation}
 E_1(z)=E_i e^{\mathrm{i}k_1z}+E_r e^{-\mathrm{i}k_1z},
\end{equation}
ndërsa në mjedisin e dytë $E_2(z)=E_t e^{\mathrm{i}k_2z}$. Vazhdimësia e komponentëve tangjencialë të fushave elektrike dhe magnetike në $z=0$ jep
\begin{equation}
 1+r=t,
 \qquad
 \frac{1-r}{Z_1}=\frac{t}{Z_2},
\end{equation}
ku $r=E_r/E_i$, $t=E_t/E_i$ dhe $Z_j$ është impedanca e mjedisit. Zgjidhja është
\begin{equation}
 r=\frac{Z_2-Z_1}{Z_2+Z_1},
 \qquad
 t=\frac{2Z_2}{Z_2+Z_1}.
\end{equation}
Koeficienti energjetik i reflektimit është $R=|r|^2$; transmetimi duhet të përfshijë raportin e flukseve dhe nuk është gjithmonë thjesht $|t|^2$.

Për plazmën e ftohtë pa përplasje, $\varepsilon_r=1-\omega_p^2/\omega^2$. Kur $\omega>\omega_p$, numri valor është real dhe vala transmetohet pjesërisht. Kur $\omega<\omega_p$, $k$ është imagjinar: fusha depërton vetëm me një gjatësi karakteristike dhe nuk transporton fuqi në thellësi në modelin ideal. Në plazmë me përplasje, permitiviteti bëhet kompleks; pjesa imagjinare përfaqëson humbjen dhe ndan absorbimin nga reflektimi.

\subsubsection{Ndërfaqja në rrjet dhe gabimi numerik}
Në diferencat e fundme, koeficienti material mund të ndryshojë ndërmjet nyjeve. Vendosja e ndërfaqes pikërisht në një nyje ose gjysmëhapi ndikon në rendin lokal. Një diskretizim konservativ i ekuacionit
\begin{equation}
 u_{tt}=\partial_x\left(c(x)^2u_x\right)
\end{equation}
përdor flukse në gjysmënyje,
\begin{equation}
 \left[\partial_x(c^2u_x)\right]_j
 \approx\frac{c_{j+1/2}^2(u_{j+1}-u_j)-c_{j-1/2}^2(u_j-u_{j-1})}{\Delta x^2}.
\end{equation}
Mesatarja e koeficientit në ndërfaqe duhet zgjedhur nga kushtet e vazhdimësisë; në probleme fluksi, mesatarja harmonike shpesh ruan më mirë fizikën se mesatarja aritmetike.

Për të matur reflektimin numerik, sinjali regjistrohet në një sondë para ndërfaqes. Simulimi referencë pa ndërfaqe jep pulsin rënës; diferenca izolon pulsin e reflektuar. Dritaret kohore duhet të ndajnë pulset dhe kufiri i majtë nuk duhet të reflektojë para përfundimit të matjes. Koeficienti mund të nxirret nga amplituda, energjia ose transformimi Fourier, por përkufizimi duhet të përputhet me formulën teorike.

\subsubsection{Absorbimi kufitar}
Një domen i fundëm duhet të imitojë hapësirën e pafundme. Kushti $u_t+cu_x=0$ është i saktë për valë njëdrejtimëshe në një përmasë dhe incidencë normale, por jo për të gjitha frekuencat dhe këndet. Një shtresë amortizuese shton termin $-\sigma(x)u_t$, me $\sigma$ që rritet gradualisht drejt kufirit. Rritja e menjëhershme e $\sigma$ krijon vetë një ndërfaqe dhe reflektim.

Shtresa e përputhur në mënyrë të përsosur (PML) projektohet që impedanca të mbetet e përputhur në hyrje dhe vala të shuhet brenda shtresës. Implementimi i saj është më i ndërlikuar, por shërben si shembull se kushti kufitar është pjesë e modelit numerik. Në laborator mund të krahasohen kufiri i fiksuar, kushti dalës i rendit të parë dhe shtresa amortizuese duke matur energjinë e kthyer.

\subsubsection{Bilanci diskret i energjisë}
Për skemën qendrore, shpejtësia diskrete mund të vendoset në gjysmëhapa. Një energji e përshtatshme është
\begin{equation}
 E^{n+1/2}=\frac{\Delta x}{2}\sum_j
 \left(\frac{u_j^{n+1}-u_j^n}{\Delta t}\right)^2
 +\frac{c^2\Delta x}{2}\sum_j
 \left(\frac{u_{j+1}^{n+1}-u_j^{n+1}}{\Delta x}\right)
 \left(\frac{u_{j+1}^{n}-u_j^{n}}{\Delta x}\right).
\end{equation}
Për kufij të përshtatshëm dhe CFL të plotësuar, kjo madhësi mbetet e kufizuar. Monitorimi i saj zbulon paqëndrueshmërinë shumë më herët se inspektimi vizual i animacionit.
